\chapter{Observation, Measurement, and the Making of Data}
\label{ch:observation}

\begin{chapterguide}
This chapter examines how observations become data. It argues that theory
ladenness does not make observation arbitrary. Instead, objectivity depends on
a traceable measurement path, calibrated instruments, controlled sources of
error, and the convergence of distinct ways of making the target answer.
\end{chapterguide}

A datum is not a small piece of reality waiting to be collected. It is the
result of an interaction between a target, an instrument, a procedure, a
record, and an interpretive model. The fact that data are made does not mean
that they are made up. A footprint is produced by a foot and a surface, but
the scientific significance of the footprint depends on a method of reading
it.

This point matters because scientific disputes often focus on whether an
observation is theory-laden. The correct answer is yes: observations rely on
concepts and instruments. But the conclusion that observations are therefore
subjective does not follow. A map is concept-laden and can still be constrained
by terrain. The relevant question is whether the concepts and instruments are
arranged so that the result can resist convenient reinterpretation.

\section{Theory-ladenness without observational chaos}

Norwood Hanson and later philosophers emphasized that what a scientist sees
depends on training and expectation. A radiologist and a novice may look at
the same image and notice different patterns. A physicist sees a detector event
as evidence of a particle only within a theoretical and instrumental context.
Wilfrid Sellars argued that the \enquote{given} is not a self-interpreting
foundation of knowledge \citep{sellars1963}.

The strongest lesson is semantic and inferential. There is no observation
sentence whose meaning is entirely independent of all concepts. The novice and
the expert do not occupy identical epistemic positions. But the existence of
different interpretations also creates a testable difference. An expert's
interpretation can be assessed by calibration, successful prediction,
independent readings, and the ability to distinguish real signals from
artifacts.

Theory-ladenness has two directions. Theory guides attention and helps design
instruments. Evidence also disciplines theory by exposing predictions that
fail, measurements that do not replicate, and interventions that produce
unexpected effects. The relation is therefore reciprocal rather than one-way.

\section{The measurement path}

To evaluate a measurement, \RCR\ asks for a path from target to claim:

\begin{enumerate}
\item \textbf{Target:} What property, event, or process is being measured?
\item \textbf{Interaction:} How does the target affect the instrument or
observer?
\item \textbf{Calibration:} How is the instrument's response related to a
standard or independently measured quantity?
\item \textbf{Transformation:} What mathematical or computational operations
turn the raw signal into a reported value?
\item \textbf{Error model:} Which systematic and random errors could distort
the result?
\item \textbf{Interpretation:} What claim is the result taken to support, and
what background assumptions connect it to the target?
\end{enumerate}

This path is not a demand that every measurement be philosophically
reconstructed from scratch. It is a practical test of where objectivity could
fail. A value can be precise but invalid if it measures the wrong construct. A
survey can be representative but biased by nonresponse. A detector can be
sensitive but unable to distinguish signal from background. A psychological
scale can be reliable across repeated administrations while failing to measure
the concept its label suggests.

\begin{figure}[ht]
\centering
\begin{tikzpicture}[
  node distance=6mm and 8mm,
  box/.style={draw=linegray,rounded corners=2pt,fill=pale,align=center,minimum width=2.55cm,minimum height=0.72cm,font=\small},
  arrow/.style={-{Latex[length=2mm]},draw=teal,thick}
]
\node[box] (target) {Target property\\or event};
\node[box,right=of target] (interaction) {Physical or social\\interaction};
\node[box,right=of interaction] (instrument) {Instrument and\\calibration};
\node[box,below=of instrument] (transform) {Data processing\\and error model};
\node[box,left=of transform] (result) {Reported\\result};
\draw[arrow] (target) -- (interaction);
\draw[arrow] (interaction) -- (instrument);
\draw[arrow] (instrument) -- (transform);
\draw[arrow] (transform) -- (result);
\end{tikzpicture}
\caption{A measurement path. Every arrow is a potential site of error and
also a site where independent validation can be added.}
\label{fig:measurement-path}
\end{figure}

\section{Operationalism and construct validity}

Operationalism insists that a concept should be defined by an operation. The
view has a useful corrective: a definition that cannot guide any procedure is
not yet an empirical concept. But one concept may have several legitimate
operations. Temperature can be measured through expansion, resistance,
radiation, or thermodynamic relations. Intelligence, poverty, and social trust
also admit multiple operationalizations.

An operation is therefore not identical to a definition. It is an instrument
for producing evidence about a construct under conditions. Convergent
measurement asks whether different operations agree when they should.
Discriminant measurement asks whether they remain distinct from neighboring
constructs. Known-group tests ask whether a measure behaves differently in
groups that the theory says should differ. These tests do not eliminate
interpretation, but they distribute it across independent routes.

\section{From data to phenomena}

James Bogen and James Woodward distinguished data from phenomena
\citep{bogenwoodward1988}. Data are local, often noisy records: a set of
instrument readings, images, or observations. Phenomena are stable patterns
that data reveal through analysis, replication, and background knowledge. The
distinction explains why a single observation can be weak evidence while a
robust phenomenon can support a theory.

For example, a particular climate station reading may be affected by
instrument placement, urbanization, or missing values. A long-run warming
pattern is not simply one more reading; it is a phenomenon supported by many
records, corrections, physical mechanisms, and independent indicators. A
single detector event may be ambiguous, while a distribution of events with a
predicted energy relation is a phenomenon.

The distinction also protects against two errors. Data are not automatically
theory-free facts, but phenomena are not inventions of a theory. A phenomenon
is a stabilized pattern that remains identifiable when the routes to it are
varied and criticism is allowed to operate.

\section{Error, uncertainty, and precision}

Precision concerns the spread of repeated measurements. Accuracy concerns
closeness to the target value. A measurement can be precise but systematically
wrong. A model can estimate a tiny standard error because it ignores a major
source of uncertainty. A survey can report many decimal places while its
conceptual validity is poor.

\RCR\ distinguishes three forms of uncertainty.

\begin{itemize}
\item \textbf{Aleatory uncertainty} concerns variation in a process or outcome
that is treated as stochastic for the purpose of analysis.
\item \textbf{Epistemic uncertainty} concerns limited knowledge of parameters,
models, initial conditions, or measurement error.
\item \textbf{Scope uncertainty} concerns whether a result transfers to a
different population, scale, time, or institutional setting.
\end{itemize}

The distinction is not merely terminological. More observations may reduce
sampling uncertainty but do little to correct a biased measurement. A larger
sample may make a misleading operational definition more precise. A model
ensemble may explore parameter uncertainty while leaving structural uncertainty
untouched.

\section{Instrumental realism}

Ian Hacking's emphasis on experimentation suggests a useful form of realism:
instruments can create a stable relation between a target and an effect. But
instrumental success must be unpacked. A detector's output can be used to
control a beam, diagnose a condition, or produce a new phenomenon. The stronger
the independence between the instrument's construction and the claim being
tested, the stronger the resulting support.

One should ask:

\begin{enumerate}
\item Can the result be reproduced with a different instrument principle?
\item Can the calibration be checked against an independent standard?
\item Does the target support interventions that were not used to define it?
\item Do failed attempts reveal a boundary rather than disappear through
post hoc interpretation?
\end{enumerate}

This standard is demanding but not impossible. It is common in metrology,
particle physics, medicine, and engineering. In social research, the analogous
tests may involve alternative measures, preregistered coding, blinded
assessment, administrative records, or natural experiments.

\section{Observation in the social sciences}

Social categories present a special issue because measurement can change the
thing measured. A census category can become an administrative identity. A
test score can alter self-perception. A policy indicator can become a target
that organizations optimize. This does not make the category unreal, but it
means that the measurement relation is reflexive.

\begin{principlebox}
\textbf{Reflexive measurement principle.} When the act of classification can
change behaviour, the measurement path must include feedback from the
classification to the target. Objectivity requires reporting that feedback,
not pretending that the target is static.
\end{principlebox}

Reflexive measurement is one reason why social science needs both causal
models and institutional analysis. A variable can be a real social fact
because people and organizations act on it, while its meaning and effects
depend on rules that can change.

\section*{Chapter summary}

Observation is conceptually mediated but not arbitrary. A measurement becomes
objective to the extent that its path from target to reported result is
traceable, calibrated, error-aware, and open to independent checks. Data are
local records; phenomena are stabilized patterns that survive analysis and
replication. Precision does not guarantee validity, and social measurements
require attention to feedback. Theory and evidence constrain one another
through carefully designed measurement paths.

\begin{discussion}
\begin{enumerate}
\item Which is more dangerous in a measurement: random noise, systematic bias,
or a construct that is poorly defined?
\item Can two measurements of the same construct legitimately disagree?
\item How should a researcher report a result when the target changes in
response to being measured?
\end{enumerate}
\end{discussion}

\begin{furtherreading}
See Hacking (1983) on experiment and intervention, Bogen and Woodward (1988)
on data and phenomena, Sellars (1963) on the critique of the given, and
Frigg and Hartmann (2006) for models and measurement in scientific practice.
\end{furtherreading}

